\label{chap:83-06}

\section{Mesure projective et conservation évolutive de l'unité}
\label{p12-83-c6-s1:chap:medida-proyectiva}

La topologie des cylindres permet de parler de proximité ; une théorie des chemins
a en outre besoin de poids. La conservation évolutive de l'unité acquiert ici
une forme probabiliste exacte : chaque cylindre répartit son poids entre ses
extensions et la somme totale demeure égale à un. L'unité ne reste pas fixe dans
une cellule ; elle est transportée à travers le raffinement.

\subsection{Poids cohérents sur les niveaux finis}

Soit \(\mu_n\) une probabilité sur l'ensemble fini
\(\mathcal S_n\). La famille est \emph{projectivement cohérente} lorsque
\begin{equation}
 \mu_n(a)
 =\sum_{\substack{b\in\mathcal S_{n+1}\\
                   \rho_{n+1,n}(b)=a}}
   \mu_{n+1}(b)
 \qquad(a\in\mathcal S_n).
 \label{p12-83-c6-s1:eq:consistencia-pesos}
\end{equation}
Cette égalité exprime la conservation locale de l'unité : le poids d'une
histoire tronquée est la somme des poids de tous ses enfants.

\begin{theorem}[Mesure généalogique sur les cylindres]
\label{p12-83-c6-s1:thm:medida-genealogica}
Toute famille de probabilités finies qui satisfait
\eqref{p12-83-c6-s1:eq:consistencia-pesos} détermine une unique probabilité borélienne
\(\mu_\infty\) sur \(\Omega_\infty\) telle que
\begin{equation}
 \mu_\infty([a]_n)=\mu_n(a).
 \label{p12-83-c6-s1:eq:medida-cilindro}
\end{equation}
\end{theorem}

\begin{proof}
\[
 \mathscr A
 :=\{\text{réunions finies de cylindres}\}
\]
est une algèbre : l'intersection de deux cylindres est vide ou est un cylindre du
niveau le plus profond. Tout élément de \(\mathscr A\) peut être raffiné en une
réunion disjointe de cylindres d'un même niveau. Si
\(A=\bigsqcup_{a\in I}[a]_n\), nous posons
\[
 \mu_0(A)=\sum_{a\in I}\mu_n(a).
\]
Itérer \eqref{p12-83-c6-s1:eq:consistencia-pesos} démontre que la valeur ne dépend pas du
niveau commun choisi.

Si un cylindre \([a]_n\) est vide, la ramification finie et le lemme de König
\cite{Koenig1936}
impliquent qu'il existe une profondeur sans descendants de \(a\). La
cohérence itérée donne alors \(\mu_n(a)=0\). C'est pourquoi \(\mu_0\) est bien
définie même pour les représentations contenant des cylindres vides et est
finiment additive.

Les cylindres sont ouverts et fermés. Si une suite décroissante
d'éléments de \(\mathscr A\) avait une intersection vide, la compacité
impliquerait que l'un d'eux est déjà vide ; c'est pourquoi \(\mu_0\) est continue en
l'ensemble vide et constitue une prémesure dénombrablement additive. Le théorème d'extension de
Carathéodory ---de manière équivalente, l'extension de Kolmogorov pour ce système
projectif fini--- la prolonge à la \(\sigma\)-algèbre engendrée par les
cylindres \cite{Caratheodory1914,Kolmogorov1933}. Celle-ci est la tribu borélienne de
\(\Omega_\infty\), et la finitude de la
mesure garantit l'unicité.
\end{proof}

La mesure peut provenir de transitions locales. Si
\(P_n(b\mid a)\ge0\) satisfait
\[
 \sum_{\rho_{n+1,n}(b)=a}P_n(b\mid a)=1,
\]
alors
\(\mu_{n+1}(b)=\mu_n(a)P_n(b\mid a)\) est cohérente. L'uniformité est un
choix possible, non une obligation. Les poids peuvent dépendre du feuillet,
de la retenue, de l'orientation, d'un coût d'action non négatif ou d'un registre pourvu
qu'ils préservent la somme. Par exemple,
\[
 P_n(b\mid a)
 =\frac{e^{-\beta a_n(b\mid a)}}
        {\sum_{b'\mapsto a}e^{-\beta a_n(b'\mid a)}},
 \qquad a_n\geq0,
\]
définit des poids positifs normalisés. La phase orthogonale d'action appartient au
noyau du \cref{chap:integral-genealogica} ; elle n'est pas insérée dans ces
probabilités positives.

\subsection{Désintégration de la mesure généalogique sur les fibres de même valeur archimédienne}

L'évaluation archimédienne transporte la mesure vers \(K\) :
\begin{equation}
 \nu:=\Val_*\mu_\infty,
 \qquad
 \nu(B)=\mu_\infty(\Val^{-1}(B)).
 \label{p12-83-c6-s1:eq:pushforward-arquimediano}
\end{equation}
La probabilité visible d'une région additionne le poids de toutes les généalogies qui
la publient. Lorsque deux sections possèdent la même valeur, leurs contributions
s'agrègent dans \(\nu\), tout en demeurant séparées dans \(\mu_\infty\).

\begin{proposition}[Désintégration par valeurs]
Si \(K\) est un espace borélien standard, il existe une famille de mesures
conditionnelles \(\{\mu_y\}_{y\in K}\), définie pour \(\nu\)-presque tout \(y\),
telle que \(\mu_y\) est portée par \(\mathfrak G_y=\Val^{-1}(y)\) et
\[
 \int_{\Omega_\infty}F\dd\mu_\infty
 =\int_K\left(\int_{\mathfrak G_y}F\dd\mu_y\right)\dd\nu(y)
\]
pour toute fonction intégrable \(F\).
\end{proposition}

Cette formule sépare valeur et généalogie au niveau de la mesure. La physique qui
ne dépend que de la valeur intègre d'abord toute la fibre ; une observable de mémoire
peut distinguer des distributions \(\mu_y\) différentes sur la même
coordonnée réelle.

\subsection{Conditionnement comme élagage des futurs}

Soit \(A\subseteq\Omega_\infty\) un événement observable avec
\(\mu_\infty(A)>0\). La mesure conditionnelle est
\begin{equation}
 \mu_\infty(B\mid A)
 =\frac{\mu_\infty(B\cap A)}{\mu_\infty(A)}.
 \label{p12-83-c6-s1:eq:condicionamiento-genealogico}
\end{equation}
Lorsque \(A\) est la lecture d'un registre, son intersection avec chaque cylindre
élimine les extensions incompatibles. La normalisation ne crée pas de nouvelles branches :
elle redistribue l'unité entre celles qui ont survécu.

\begin{proposition}[Mise à jour compatible]
Si \(A\) est une réunion de cylindres au niveau \(m\), alors pour tout \(n\ge m\)
la distribution conditionnelle sur \(\mathcal S_n\) s'obtient en supprimant les
états dont les troncatures n'appartiennent pas à \(A\) et en renormalisant les poids
restants. Les distributions conditionnelles continuent de satisfaire
\eqref{p12-83-c6-s1:eq:consistencia-pesos}.
\end{proposition}

\begin{proof}
L'identité est la règle de Bayes sur les ensembles finis. La somme des
poids des enfants conditionnés coïncide avec le poids conditionné du parent,
car intersection et réunion disjointe commutent.
\end{proof}

\subsection{Conservation combinatoire et complétion dimensionnelle}

La même loi apparaît dans la complétion binomiale :
\[
 1=\sum_{r=0}^{d}\binom dr
 \left(\frac9{10}\right)^{d-r}
 \left(\frac1{10}\right)^r.
\]
En dimension trois,
\[
 1000=729+243+27+1,
\]
de sorte que l'intérieur, les nouvelles faces, les arêtes extrêmes et le sommet final forment
une partition de l'unité décimale. Le passage
\(729+270=999\to729+271=1000\) distingue l'état antérieur à la clôture de
l'unité complétée. Dans le système de chemins, l'identité analogue est
\[
 \mu([a]_n)=\sum_{b\mapsto a}\mu([b]_{n+1}).
\]
L'unité évolue en se distribuant dans les strates et se retrouve en additionnant
la généalogie complète.

\begin{figure}[htbp]
\centering
\begin{tikzpicture}[>=Latex,node distance=12mm and 17mm,
  box/.style={draw=hmtblue,rounded corners,fill=hmtlightblue,align=center,
              minimum width=28mm,minimum height=8mm},
  leaf/.style={draw=hmtgreen,rounded corners,fill=hmtlightgreen,align=center,
               minimum width=22mm,minimum height=7mm}]
\node[box] (a) {$[a]_n$\\poids $\mu_n(a)$};
\node[leaf,below left=of a] (b1) {$[b_1]_{n+1}$\\$\mu_{n+1}(b_1)$};
\node[leaf,below=of a] (b2) {$[b_2]_{n+1}$\\$\mu_{n+1}(b_2)$};
\node[leaf,below right=of a] (b3) {$[b_3]_{n+1}$\\$\mu_{n+1}(b_3)$};
\draw[->] (a)--(b1); \draw[->] (a)--(b2); \draw[->] (a)--(b3);
\node[below=20mm of b2,align=center,text width=.78\textwidth] (eq)
 {$\mu_n(a)=\mu_{n+1}(b_1)+\mu_{n+1}(b_2)+\mu_{n+1}(b_3)$ :\\
  l'unité se conserve lorsque la généalogie est raffinée.};
\end{tikzpicture}
\caption{Conservation projective du poids sur les cylindres.}
\label{p12-83-c6-s1:fig:conservacion-peso-cilindros}
\end{figure}

\begin{statusbox}
La mesure projective se construit sur le système projectif HMT.
Le théorème d'extension relève des mathématiques classiques ; l'identification de ses
cylindres aux états TPK est exacte une fois les poids de transition fixés.
La sélection d'une mesure physique concrète appartient au modèle dynamique et ne
se déduit pas uniquement du recensement des branches.
\end{statusbox}

La mesure projective est caractérisée par la positivité, la conservation du
poids entre parents et enfants et le conditionnement au sein de l'événement observé. Une
réalisation expérimentale complète ajoute la préparation et le protocole
d'échantillonnage ; avec ces données, la même mesure produit des probabilités de lecture
comparables à des fréquences.
